\chapter{Model, observables, and physical interpretation}\label{sec:model}
\section{Dynamics and observables}\label{sec:observables}
Let $\Lambda_L=\{0,\ldots,L-1\}^2$ be the lattice, with stable heights $h(x)\in\{0,1,2,3\}$. We use the local rule in the project specification~\cite{handout}. Starting from an empty grid, the primary baseline chooses a site independently and uniformly, adds one grain and lets the grid relax completely. A legal elementary toppling requires $h(x)\geq4$, subtracts four grains and transfers one in each axial direction. Grains sent outside the lattice leave the system. The threshold remains four at edges and corners. The next grain is added only after stabilization, and external time counts input events.

Let $u_t(x)$ count elementary topplings at site $x$ in event $t$. The sitewise toppling-count vector $u_t$ is called the odometer. Let $X_t$ be the set of input sites. We measure
\begin{align}
 S_t&=\sum_xu_t(x), & A_t&=\sum_x\one_{\{u_t(x)>0\}},\\
 D_t&=\sum_x[4-d_G(x)]u_t(x),&
 \ell_{X_t}&=\begin{cases}1+\displaystyle\max_{u_t(x)>0}\min_{s\in X_t}d_G^{\rm geo}(x,s),&S_t>0,\\0,&S_t=0.\end{cases}
 \label{eq:observables}
\end{align}
Here $d_G(x)$ is the number of internal directional neighbours and $d_G^{\rm geo}$ is graph distance. On the open square, distance is Manhattan distance. With one input site, reach is the greatest distance from that site to a toppled site, plus one. A single-site avalanche therefore has reach one. For several input sites, we measure each toppled site from its nearest source.

We calculate mean $S$ and $D$ over every event, including $S=0$. Per-grain values divide the totals by the number of grains added. The main avalanche distributions and correlations condition on $S>0$ and retain $D=0$. Positive-loss fits condition on $D>0$. We report the zero-event probabilities separately. Mean height uses total mass after every event, while height frequencies are sampled every 128 added grains after stabilization, including after zero events. For batches, elapsed input is measured in grains as well as events. The distance from a batch to a dissipative boundary is the minimum over its input sites.

\section{Scale invariance and the SOC mechanism}\label{sec:scaling}
Scale invariance means that changing scale leaves the shape of a relation unchanged, apart from an overall factor. Suppose a positive function satisfies $f(cx)=g(c)f(x)$ for every $c,x>0$. Rescaling first by $a$ and then by $b$ gives $g(ab)=g(a)g(b)$. Under continuity or suitable measurability, $g(c)=c^{-\alpha}$, giving
\begin{equation}
 f(x)=f(x_0)(x/x_0)^{-\alpha}=Cx^{-\alpha}.
\end{equation}
For an unbounded discrete power tail with $\alpha>1$, $\Prob(X\geq n)\sim C'n^{-(\alpha-1)}$, so the probability mass function (PMF) and complementary cumulative distribution function (CCDF) have different exponents. On a finite lattice, microscopic and system-size cutoffs limit the range over which scaling can hold approximately.

The separation between slow input and rapid relaxation is central to the BTW model. Local thresholds allow activity to spread, interior topplings conserve grains and the boundary removes them. These features motivate examining height, transport and avalanche distributions together. We assess power-law behaviour by fitting normalized distributions over specified ranges~\cite{clauset2009}.

\section{Boundaries, steady states, and infinite volume}\label{sec:boundaries}
If $m_t$ grains enter an event, total mass satisfies
\begin{equation}
 M_t=M_{t-1}+m_t-D_t,\qquad
 \overline D=\overline m-\frac{M_T-M_0}{T}.\label{eq:mass}
\end{equation}
Open boundaries allow grains to leave as new grains arrive. Every site connects to a sink, and stable mass is bounded by $3L^2$. The endpoint term in Eq.~\eqref{eq:mass} therefore vanishes over a long window. Average loss can approach average input while other statistics still show transients. Without dissipation, continued input to a finite system eventually exceeds the capacity of any stable configuration.

Starting from an empty infinite lattice, any finite collection of grains can stabilize by spreading outward~\cite{levinepropp2010}. Continued input in a finite source region increases the total mass and expands the affected region.

We use \emph{macroscopic steady state} to mean that selected bulk observables show no systematic drift. \emph{Statistical stationarity} is stronger: the joint probability laws remain unchanged under shifts in time. In the physical sense used here, \emph{equilibrium} has no sustained externally driven transport. Detailed balance is a common criterion in stochastic models. The relevant description of the driven sandpile is a nonequilibrium stationary state. Grains continue to enter and leave even when the height statistics settle.
